% Part:sets-functions-relations
% Chapter: size-of-sets
% Section: equinumerous-sets

\documentclass[../../../include/open-logic-section]{subfiles}

\begin{document}

\olfileid{sfr}{siz}{equ}

\olsection{Padha Cacahe}

Kita duwe gagasan intuisi babagan "gedhene" himpunan, kang lumaku kanthi
becik kanggo himpunan winates. Nanging kepriye yen himpunane tanpa wates?
Yen arep nemokake cara formal kanggo mbandhingake gedhene rong himpunan
kang gedhene \emph{sembarang}, prayoga miwiti kanthi netepake kapan
rong himpunan padha gedhene. Frege ngandharake:
\begin{quote}
  Yen sawijining pelayan arep mesthekake yen cacah peso kang ditata ing meja
  persis padha karo cacah piring, dheweke ora perlu ngetung loro-lorone.
  Cukup nyelehake siji peso ing sisih tengen saben piring, saengga saben
  peso ing meja ana ing sisih tengen sawijining piring. Kanthi mangkono,
  piring lan peso digandhengake siji karo sijine kanthi tunggal,
  lan kanthi sesambungan papan kang padha. \citep[\S70]{Frege1884}
\end{quote}
Gagasan ing pethikan iki bisa diandharake nganggo definisi formal:

\begin{defn}\ollabel{comparisondef}
  $A$ \emph{padha cacahe} karo $B$, ditulis $\cardeq{A}{B}$,
  yen lan mung yen ana !!a{bijection} $f \colon A \to B$.
\end{defn}

\begin{prop}\ollabel{equinumerosityisequi}
Padha cacahe iku relasi ekuivalensi.
\end{prop}

\begin{proof}
Kita kudu nuduhake yen relasi padha cacahe iku refleksif, simetris,
lan transitif. Upamakna $A, B$, lan $C$ himpunan.

\emph{Refleksif.} Pemetaan identitas $\Id{A} \colon A \to A$, kanthi
$\Id{A} (x) = x$ kanggo kabeh $x \in A$, iku !!a{bijection}.
Dadi $\cardeq{A}{A}$.

\emph{Simetris.} Upamakna $\cardeq{A}{B}$, yaiku ana
!!a{bijection} $f\colon A \to B$. Amarga $f$ iku !!{bijective},
inverse $f^{-1}$ ana lan uga !!{bijective}. Mula,
$f^{-1}\colon B \to A$ iku !!a{bijection}, dadi $\cardeq{B}{A}$.

\emph{Transitif.} Upamakna $\cardeq{A}{B}$ lan $\cardeq{B}{C}$,
yaiku ana !!{bijection}s $f\colon A \to B$ lan $g\colon B \to
C$. Mula komposisi $\comp{f}{g}\colon A \to C$ iku !!{bijective},
saengga $\cardeq{A}{C}$.
\end{proof}

\begin{prop}
Yen $\cardeq{A}{B}$, mula $A$ iku !!{enumerable}
yen lan mung yen $B$ uga mangkono.
\end{prop}

\begin{editorial}
Bukti ing ngisor nggunakake \olref[enm]{defn:enumerable} yen
\olref[enm]{sec} kalebu ing wacan, lan \olref[enm-alt]{defn:enumerable}
yen ora.
\end{editorial}

\begin{proof}
Upamakna $\cardeq{A}{B}$, mula ana !!{bijection} $f \colon A
\to B$, lan upamakna $A$ iku !!{enumerable}.
\oliflabeldef{sfr:siz:enm:defn:enumerable}{
  Dadi $A = \emptyset$ utawa ana fungsi !!a{surjective}
  $g\colon \PosInt \to A$. Yen $A = \emptyset$, mula $B = \emptyset$
  uga; yen ora, ana !!a{element}~$y \in B$ nanging ora ana $x \in
  A$ kanthi $g(x) = y$. Nanging yen $g\colon \PosInt \to A$
  iku !!{surjective}, mula $\comp{g}{f} \colon \PosInt \to B$
  iku !!{surjective}. Kanggo mriksa iki, upamakna $y \in B$. Amarga $f$
  iku !!{surjective}, ana $x \in A$ saengga $f(x) = y$. Amarga
  $g$ iku !!{surjective}, ana $n \in \PosInt$ saengga $g(n) =
  x$. Mula,
\[
(\comp{g}{f})(n) = f(g(n)) = f(x) = y
\]
dadi $\comp{g}{f}$ iku !!{surjective}. Kita oleh yen $\comp{g}{f}$
iku enumerasi saka~$B$, mula $B$~iku !!{enumerable}.}
{
  Dadi $A  = \emptyset$ utawa ana !!a{bijection}~$g$ kang
  jangkauane $A$ lan domaine $\Nat$ utawa rerangken wiwitan
  wilangan natural. Yen $A = \emptyset$, mula $B =
  \emptyset$ uga; yen ora, ana~$y \in B$ nanging ora ana $x
  \in A$ saengga $g(x) = y$. Mula upamakna ana
  !!{bijection}~$g$ mau. Dadi $\comp{g}{f}$ iku !!a{bijection} kanthi
  jangkauan~$B$ lan domain kang padha karo domaine~$g$, yaiku $\Nat$ utawa
  perangan wiwitane, saengga $B$ iku !!{enumerable}.}

Yen $B$ iku !!{enumerable}, kita bisa nuduhake yen $A$ iku !!{enumerable}
kanthi mbaleni argumen nganggo !!{bijection} $f^{-1}\colon B \to A$
kanggo ngganti~$f$.
\end{proof}

\begin{prob}
Tuduhna yen $\cardeq{A}{C}$ lan $\cardeq{B}{D}$, sarta $A \cap B =
C \cap D = \emptyset$, mula $\cardeq{A \cup B}{C \cup D}$.
\end{prob}

\begin{prob}
Tuduhna yen $A$ tanpa wates lan !!{enumerable}, mula
$\cardeq{A}{\Nat}$.
\end{prob}

\end{document}
